\documentclass[11pt]{article}
% ======================================================================
%  examstyle.tex — EPFL-style exam layout for the weekly Time Series exams
%  Shared by every Exam_Week_NN(.tex / _Solutions.tex). Compiles with pdflatex.
% ======================================================================
\usepackage[utf8]{inputenc}
\usepackage[T1]{fontenc}
\usepackage[english]{babel}
\usepackage[a4paper,margin=2.4cm]{geometry}
\usepackage{amsmath,amssymb,amsthm,mathtools}
\usepackage{bm}
\usepackage{enumitem}
\usepackage{array}

\setlength{\parindent}{0pt}
\setlength{\parskip}{0.5em}
\emergencystretch=3em
\allowdisplaybreaks[1]

% ---------- math macros (same names as the rest of the project) ----------
\newcommand{\E}{\mathbb{E}}
\DeclareMathOperator{\Var}{Var}
\DeclareMathOperator{\Cov}{Cov}
\DeclareMathOperator{\Corr}{Corr}
\DeclareMathOperator*{\argmin}{arg\,min}
\DeclareMathOperator{\MSE}{MSE}
\DeclareMathOperator{\trace}{tr}
\newcommand{\Reals}{\mathbb{R}}
\newcommand{\Ints}{\mathbb{Z}}
\newcommand{\Comp}{\mathbb{C}}
\newcommand{\Nat}{\mathbb{N}}
\newcommand{\Prob}{\mathbb{P}}
\newcommand{\Tr}{^{\mathsf{T}}}
\newcommand{\conj}[1]{\overline{#1}}
\newcommand{\abs}[1]{\left\lvert #1 \right\rvert}
\newcommand{\norm}[1]{\left\lVert #1 \right\rVert}
\newcommand{\ind}{\mathbf{1}}
\newcommand{\eps}{\varepsilon}
\newcommand{\diff}{\nabla}
\newcommand{\acvs}{\gamma}
\newcommand{\acf}{\rho}
\newcommand{\pacf}{\alpha}
\newcommand{\sdf}{\mathcal{S}}
\newcommand{\Bsh}{B}
\newcommand{\iid}{\overset{\text{iid}}{\sim}}

\setlist[enumerate]{leftmargin=2em,topsep=2pt,itemsep=4pt}
\setlist[itemize]{leftmargin=1.6em,topsep=2pt,itemsep=2pt}

\newcounter{exo}

% ---------- exam cover header :  \examheader{exam title}{duration} ----------
\newcommand{\examheader}[2]{%
  \thispagestyle{empty}
  \begin{center}
    {\Large\bfseries Time Series}\\[3pt]
    Spring Semester 2026, EPFL\\
    \'Ecole Polytechnique F\'ed\'erale de Lausanne\\
    Prof.\ Dr.\ Sofia C.\ Olhede\\[7pt]
    \hrule height 0.8pt
    \vspace{9pt}
    {\large\bfseries Practice Exam --- #1}\\[4pt]
    {\normalsize Duration: #2 \quad$\cdot$\quad No calculator \quad$\cdot$\quad Answer on paper}\\
    \vspace{8pt}
    \hrule height 0.8pt
  \end{center}
  \vspace{14pt}
  \begin{tabular}{@{}p{8.2cm}@{}p{8cm}@{}}
    Name:\enspace\hrulefill & Surname:\enspace\hrulefill
  \end{tabular}\\[14pt]
  SCIPER (Student Card Number):\enspace\hrulefill
  \vspace{16pt}
}

% ---------- points table (fixed: 4 problems) ----------
\newcommand{\pointstable}{%
  \begin{center}
  \renewcommand{\arraystretch}{1.5}
  \begin{tabular}{|p{5cm}|p{4cm}|}
    \hline
    \textbf{Exercise} & \textbf{Points}\\\hline
    1 & \\\hline
    2 & \\\hline
    3 & \\\hline
    4 & \\\hline
    \textbf{Total} & \\\hline
  \end{tabular}
  \end{center}
}

% ---------- problem heading (exam: one problem per page) ----------
\newcommand{\problem}[1]{%
  \clearpage
  \stepcounter{exo}%
  \begin{center}\textbf{\large Problem \theexo\ (#1 points)}\end{center}
  \vspace{4pt}
}
% ---------- answer space: fill the statement page + one blank answer page ----------
% Put \answerpage at the END of each problem so every exercise gets its own page(s).
\newcommand{\answerpage}{\par\vfill\newpage\null\newpage}

% ---------- solutions document ----------
\newcommand{\solheader}[1]{%
  \begin{center}
    {\Large\bfseries Time Series --- Solutions}\\[3pt]
    {\large Practice Exam --- #1}\\[2pt]
    EPFL, MATH-342 \quad$\cdot$\quad Prof.\ S.\ C.\ Olhede\\[5pt]
    \hrule height 0.8pt
  \end{center}
  \vspace{10pt}
}
\newcommand{\solproblem}[1]{%
  \bigskip
  \stepcounter{exo}%
  {\large\bfseries Problem \theexo\ (#1 points)}\par\nobreak\vspace{2pt}
}
% Rappel de l'énoncé avant la solution (dans les corrigés)
\newcommand{\statementstart}{\par\vspace{1pt}{\bfseries\itshape Statement.}\par\nobreak\begingroup\small\itshape}
\newcommand{\statementend}{\par\endgroup\vspace{5pt}\hrule\vspace{6pt}{\bfseries Solution.}\par\nobreak\vspace{2pt}}

\begin{document}

\solheader{Week 8: Spectral Estimation --- Periodogram \& Tapering}

% ---------------------------------------------------------------
\solproblem{5}
Write $Y_t=X_t-\bar X$ throughout.

\textbf{(a)} Substitute the definition of $\hat\acvs_\tau$ into the periodogram. Only lags $\abs\tau\le n-1$ contribute, and for each such $\tau$ the inner sum runs over $t=1,\dots,n-\abs\tau$:
\[
\hat\sdf^{(p)}(f)=\sum_{\tau=-(n-1)}^{n-1}\frac1n\sum_{t=1}^{n-\abs\tau} Y_t\,Y_{t+\abs\tau}\,e^{-2\pi i\tau f}.
\]
The pairs $(t,t+\abs\tau)$ with $1\le t$ and $t+\abs\tau\le n$ are exactly the pairs of indices $(t,s)$ with $s=t+\abs\tau\in T$, $t\in T$. As $\tau$ ranges over $-(n-1),\dots,n-1$ and $t$ over the admissible range, the pair $(t,s)$ ranges over \emph{all} of $T\times T$ once. Setting $s=t+\tau$ for the $\tau\ge0$ part and $s=t-\abs\tau$ for the $\tau<0$ part, we may rewrite the double sum as the free double sum over $(s,t)\in T\times T$ with $e^{-2\pi i\tau f}=e^{-2\pi i(s-t)f}$:
\[
\hat\sdf^{(p)}(f)
=\frac1n\sum_{t=1}^{n}\sum_{s=1}^{n} Y_t\,Y_s\,e^{-2\pi i(s-t)f}
=\frac1n\sum_{t=1}^{n}\sum_{s=1}^{n} Y_t\,Y_s\,e^{-2\pi i s f}\,e^{2\pi i t f}.
\]
The double sum factorises:
\[
\hat\sdf^{(p)}(f)
=\Bigl(\sqrt{\tfrac1n}\sum_{t=1}^{n} Y_t\,e^{2\pi i t f}\Bigr)
 \Bigl(\sqrt{\tfrac1n}\sum_{s=1}^{n} Y_s\,e^{-2\pi i s f}\Bigr)
=\conj{J(f)}\,J(f)=\abs{J(f)}^2,
\]
since $Y_t\in\Reals$ makes $\sqrt{1/n}\sum_t Y_t e^{2\pi i tf}=\conj{J(f)}$. This is the claimed identity.

\textbf{(b)} \emph{Non-negativity.} $\hat\sdf^{(p)}(f)=\abs{J(f)}^2\ge0$ for every $f$, being the squared modulus of a complex number. (No SDF estimate can be negative; the representation makes this automatic.)

\emph{Evenness.} Because the $Y_t$ are real,
\[
J(-f)=\sqrt{\tfrac1n}\sum_t Y_t\,e^{2\pi i t f}=\conj{J(f)},
\]
hence $\hat\sdf^{(p)}(-f)=\abs{J(-f)}^2=\abs{\conj{J(f)}}^2=\abs{J(f)}^2=\hat\sdf^{(p)}(f)$. So $\hat\sdf^{(p)}$ is even.

\textbf{(c)} At $f=0$ the exponential is $1$, so
\[
J(0)=\sqrt{\tfrac1n}\sum_{t=1}^{n} Y_t=\sqrt{\tfrac1n}\sum_{t=1}^{n}(X_t-\bar X)
=\sqrt{\tfrac1n}\Bigl(\sum_{t=1}^n X_t-n\bar X\Bigr)=0,
\]
because $\sum_t X_t=n\bar X$ by definition of $\bar X$. Therefore $\hat\sdf^{(p)}(0)=\abs{J(0)}^2=0$: removing the sample mean annihilates the zero-frequency ordinate.

% ---------------------------------------------------------------
\solproblem{5}
Here $\mu=0$, so $Y_t=X_t$ and $\E[Y_tY_s]=\acvs_{t-s}$.

\textbf{(a)} Take expectations in the factorised form of the periodogram (Problem 1(a)); expectation passes through the finite double sum:
\[
\E\bigl[\hat\sdf^{(p)}(f)\bigr]
=\frac1n\sum_{t=1}^{n}\sum_{s=1}^{n}\E[Y_tY_s]\,e^{-2\pi i f(t-s)}
=\frac1n\sum_{t=1}^{n}\sum_{s=1}^{n}\acvs_{t-s}\,e^{-2\pi i f(t-s)}.
\]
Collect terms by the lag $\tau=t-s$. For fixed $\tau$ with $\abs\tau\le n-1$, the number of pairs $(t,s)\in T\times T$ with $t-s=\tau$ is $n-\abs\tau$ (and there are none for $\abs\tau\ge n$). Hence
\[
\E\bigl[\hat\sdf^{(p)}(f)\bigr]
=\frac1n\sum_{\tau=-(n-1)}^{n-1}(n-\abs\tau)\,\acvs_\tau\,e^{-2\pi i\tau f}
=\sum_{\tau=-(n-1)}^{n-1}\Bigl(1-\frac{\abs\tau}{n}\Bigr)\acvs_\tau\,e^{-2\pi i\tau f},
\]
which is the stated expression with the triangular weight $w_\tau=1-\abs\tau/n$. (Equivalently, $\E[\hat\acvs_\tau]=(1-\abs\tau/n)\acvs_\tau$, then transform term by term.)

\textbf{(b)} With $d_t=\ind_T(t)$ (so $d_t=1$ for $t\in\{1,\dots,n\}$ and $0$ otherwise), the autocorrelation of the window at lag $\tau$ is
\[
\sum_{t\in\Ints} d_t\,d_{t+\abs\tau}=\#\{t: t\in T,\ t+\abs\tau\in T\}=n-\abs\tau\quad(\abs\tau\le n-1),
\]
so indeed $w_\tau=\frac1n\sum_t d_t d_{t+\abs\tau}$. Write $w=\frac1n(\tilde d * d)$ where $\tilde d_t=d_{-t}$; equivalently $w_\tau=\frac1n(\tilde d * \bar d)_\tau$ with $\bar d=d$ (the window is real). By the discrete convolution theorem, the Fourier transform of a convolution is the product of the transforms, and the transform of $\tilde d$ is $\conj{D}$ where $D(f)=\sum_{t} d_t e^{-2\pi i t f}=\sum_{t=1}^{n}e^{-2\pi i t f}$. Hence the transform of $w$ is
\[
W(f)=\frac1n\,\conj{D(f)}\,D(f)=\frac1n\,\abs{D(f)}^2=\mathcal F_n(f),
\]
the Fej\'er kernel. Now the expected periodogram is the transform of the product $w_\tau\acvs_\tau$; by the (dual) convolution theorem the transform of a product is the convolution of the transforms, $W * \sdf$:
\[
\E\bigl[\hat\sdf^{(p)}(f)\bigr]
=\sum_\tau w_\tau\acvs_\tau e^{-2\pi i\tau f}
=\int_{-1/2}^{1/2}\sdf(f')\,\mathcal F_n(f-f')\,df'.
\]
So the expected periodogram is the true SDF convolved with (blurred by) the Fej\'er kernel.

\textbf{(c)} Take $\sdf\equiv1$ in the convolution identity of (b): the corresponding ACVS is $\acvs_0=1$, $\acvs_\tau=0$ otherwise, so $\sum_\tau w_\tau\acvs_\tau e^{-2\pi i\tau f}=w_0=1$, while the right-hand side is $\int_{-1/2}^{1/2}\mathcal F_n(f-f')\,df'=\int_{-1/2}^{1/2}\mathcal F_n(u)\,du$ (periodicity, substitute $u=f-f'$). Equating,
\[
\int_{-1/2}^{1/2}\mathcal F_n(f)\,df=1.
\]
(Direct check: $\int_{-1/2}^{1/2}\mathcal F_n=\frac1n\int\abs{D}^2=\frac1n\sum_{t,s}\int_{-1/2}^{1/2}e^{-2\pi i(t-s)f}df=\frac1n\sum_{t,s}\ind_{\{t=s\}}=\frac1n\cdot n=1$.) Since $\mathcal F_n\ge0$ and integrates to $1$, it is a probability density on $[-\tfrac12,\tfrac12]$; as $n\to\infty$ it concentrates all its mass at $f=0$ (an approximate identity), so the convolution $\sdf * \mathcal F_n\to\sdf$ at every continuity point. Hence $\E[\hat\sdf^{(p)}(f)]\to\sdf(f)$: the periodogram is \emph{asymptotically} unbiased.

% ---------------------------------------------------------------
\solproblem{5}
Here $\mu$ is removed, $Y_t=X_t-\mu$, and we take $\E[Y_tY_s]=\acvs_{t-s}$.

\textbf{(a)} Expand the squared modulus and take expectations:
\[
\hat\sdf^{(p)}_h(f)=\abs{J_h(f)}^2
=\sum_{t=1}^{n}\sum_{s=1}^{n} h_t h_s\,Y_t Y_s\,e^{-2\pi i f(t-s)},
\]
\[
\E\bigl[\hat\sdf^{(p)}_h(f)\bigr]
=\sum_{t=1}^{n}\sum_{s=1}^{n} h_t h_s\,\acvs_{t-s}\,e^{-2\pi i f(t-s)}.
\]
Collect by $\tau=t-s$ (extend $h_t=0$ off $T$ so the $t$-sum is over $\Ints$):
\[
\E\bigl[\hat\sdf^{(p)}_h(f)\bigr]
=\sum_{\tau\in\Ints}\Bigl(\sum_{t\in\Ints} h_t\,h_{t-\tau}\Bigr)\acvs_\tau\,e^{-2\pi i\tau f}
=\sum_{\tau\in\Ints} w_\tau\,\acvs_\tau\,e^{-2\pi i\tau f},
\qquad
w_\tau=\sum_{t\in\Ints} h_t\,h_{t+\tau},
\]
where in the last step we re-indexed ($t\mapsto t+\tau$) and used that $w_{-\tau}=w_\tau$, so $\sum_t h_t h_{t-\tau}=\sum_t h_t h_{t+\tau}$. Now identify $w$ as the autocorrelation of the taper: with $\tilde h_t=h_{-t}$ and $\bar h_t=h_t$,
\[
(\tilde h * \bar h)_\tau=\sum_{s\in\Ints}\tilde h_{\tau-s}\,\bar h_s=\sum_{s}h_{s-\tau}h_s=w_{-\tau}=w_\tau,
\]
so $w=\tilde h * \bar h$. By the discrete convolution theorem its transform is the product of the transforms; the transform of $\tilde h$ is $\conj{H}$ and that of $\bar h$ is $H$, so
\[
W(f)=\conj{H(f)}\,H(f)=\abs{H(f)}^2 .
\]
Finally the expected tapered periodogram is the transform of the product $w_\tau\acvs_\tau$, hence the convolution of the transforms:
\[
\E\bigl[\hat\sdf^{(p)}_h(f)\bigr]=\int_{-1/2}^{1/2}\sdf(f')\,\abs{H(f-f')}^2\,df' .
\]

\textbf{(b)} For white noise $\sdf(f')\equiv\sigma^2$ is constant, so it comes out of the integral; using periodicity and the substitution $u=f-f'$,
\[
\E\bigl[\hat\sdf^{(p)}_h(f)\bigr]
=\sigma^2\int_{-1/2}^{1/2}\abs{H(f-f')}^2\,df'
=\sigma^2\int_{-1/2}^{1/2}\abs{H(u)}^2\,du .
\]
By Parseval for the discrete-time transform,
\[
\int_{-1/2}^{1/2}\abs{H(u)}^2\,du
=\int_{-1/2}^{1/2}\sum_{t}\sum_{s}h_t h_s\,e^{-2\pi i(t-s)u}\,du
=\sum_{t}\sum_{s}h_t h_s\,\ind_{\{t=s\}}
=\sum_{t} h_t^2=\norm{h}_2^2,
\]
using $\int_{-1/2}^{1/2}e^{-2\pi i\tau u}\,du=\ind_{\{\tau=0\}}$. Therefore
\[
\E\bigl[\hat\sdf^{(p)}_h(f)\bigr]=\sigma^2\,\norm{h}_2^2 .
\]
The true SDF of white noise is $\sdf(f)=\sigma^2$ for all $f$; the estimator is exactly unbiased, $\E[\hat\sdf^{(p)}_h(f)]=\sigma^2=\sdf(f)$ for \emph{every} $n$ and \emph{every} $f$, precisely when the normalisation $\norm{h}_2^2=1$ holds. (Any other normalisation would scale the estimate by the constant $\norm{h}_2^2$, biasing it; this is exactly why a taper is defined with unit energy.)

\textbf{(c)} The multitaper estimator is the average $\hat\sdf^{(mt)}(f)=\frac1K\sum_{k=1}^{K}\hat\sdf^{(p)}_{h_k}(f)$. By bilinearity of the variance,
\[
\Var\bigl(\hat\sdf^{(mt)}(f)\bigr)
=\frac{1}{K^2}\sum_{k=1}^{K}\sum_{k'=1}^{K}\Cov\bigl(\hat\sdf^{(p)}_{h_k}(f),\hat\sdf^{(p)}_{h_{k'}}(f)\bigr).
\]
Asymptotically the $K$ tapered periodograms are uncorrelated (orthogonal tapers), so that $\Cov(\hat\sdf^{(p)}_{h_k},\hat\sdf^{(p)}_{h_{k'}})\to\sdf(f)^2\,\delta_{k,k'}$. Only the $K$ diagonal terms survive, each $\to\sdf(f)^2$:
\[
\Var\bigl(\hat\sdf^{(mt)}(f)\bigr)\to\frac{1}{K^2}\cdot K\,\sdf(f)^2=\frac{\sdf(f)^2}{K}.
\]
Unlike the raw periodogram (whose variance $\to\sdf(f)^2\ne0$ and does \emph{not} shrink with $n$), here one may let $K\to\infty$ slowly as $n\to\infty$, driving $\Var\to0$ while the bias stays controlled; averaging independent estimates is what restores consistency.

% ---------------------------------------------------------------
\solproblem{5}
\textbf{(a)} Write $\theta_0=1$. For an $\mathrm{MA}(2)$, $X_t=\eps_t+\theta_1\eps_{t-1}+\theta_2\eps_{t-2}$, the ACVS is $\acvs_\tau=\sigma^2\sum_{j}\theta_j\theta_{j+\abs\tau}$ over the valid index range:
\[
\begin{aligned}
\acvs_0&=\sigma^2\bigl(\theta_0^2+\theta_1^2+\theta_2^2\bigr)=\sigma^2\bigl(1+\theta_1^2+\theta_2^2\bigr),\\
\acvs_1&=\sigma^2\bigl(\theta_0\theta_1+\theta_1\theta_2\bigr)=\sigma^2\bigl(\theta_1+\theta_1\theta_2\bigr),\\
\acvs_2&=\sigma^2\,\theta_0\theta_2=\sigma^2\,\theta_2,
\end{aligned}
\]
and $\acvs_\tau=0$ for $\abs\tau\ge3$ (no overlap of the taps), with $\acvs_{-\tau}=\acvs_\tau$. Since $\{\acvs_\tau\}\in\ell^1$ the SDF is the Fourier transform of the ACVS. Writing $\acvs_\tau=\sigma^2\sum_{j=0}^{2}\sum_{k=0}^{2}\theta_j\theta_k\ind_{\{\tau=k-j\}}$,
\[
\begin{aligned}
\sdf(f)&=\sum_{\tau\in\Ints}\acvs_\tau e^{-2\pi i\tau f}
=\sigma^2\sum_{j=0}^{2}\sum_{k=0}^{2}\theta_j\theta_k\,e^{-2\pi i(k-j)f}\\
&=\sigma^2\Bigl(\sum_{k=0}^{2}\theta_k e^{-2\pi i k f}\Bigr)\conj{\Bigl(\sum_{j=0}^{2}\theta_j e^{-2\pi i j f}\Bigr)}
=\sigma^2\,\Bigl|\sum_{k=0}^{2}\theta_k e^{-2\pi i k f}\Bigr|^2\\
&=\sigma^2\,\abs{1+\theta_1 e^{-2\pi i f}+\theta_2 e^{-4\pi i f}}^2 ,
\end{aligned}
\]
the squared modulus of the MA transfer function, as claimed (and automatically $\ge0$).

\textbf{(b)} Expanding the squared modulus (or pairing $\tau$ with $-\tau$) gives the real cosine form
\[
\sdf(f)=\acvs_0+2\acvs_1\cos(2\pi f)+2\acvs_2\cos(4\pi f).
\]
With $\theta_1=1$, $\theta_2=\tfrac14$, $\sigma^2=1$:
\[
\acvs_0=1+1+\tfrac1{16}=\tfrac{33}{16},\qquad
\acvs_1=1+\tfrac14=\tfrac54,\qquad
\acvs_2=\tfrac14 .
\]
Hence
\[
\sdf(f)=\tfrac{33}{16}+\tfrac52\cos(2\pi f)+\tfrac12\cos(4\pi f).
\]
\emph{At $f=0$:} $\cos0=1$, so $\sdf(0)=\tfrac{33}{16}+\tfrac52+\tfrac12=\tfrac{33}{16}+\tfrac{48}{16}=\dfrac{81}{16}$. (Check: $\abs{1+1+\tfrac14}^2=(\tfrac94)^2=\tfrac{81}{16}$.)

\emph{At $f=\tfrac12$:} $\cos(\pi)=-1$, $\cos(2\pi)=1$, so $\sdf(\tfrac12)=\tfrac{33}{16}-\tfrac52+\tfrac12=\tfrac{33}{16}-\tfrac{32}{16}=\dfrac{1}{16}$. (Check: at $f=\tfrac12$, $e^{-2\pi i f}=-1$, $e^{-4\pi i f}=1$, so $\abs{1-1+\tfrac14}^2=\tfrac1{16}$.)

Both values are positive. This is consistent with the process being a genuine MA: by part (a) the SDF is the squared modulus $\sigma^2\abs{\Theta(e^{-2\pi i f})}^2\ge0$ at \emph{every} $f$, so non-negativity is automatic --- exactly the property a valid SDF must have.

\textbf{(c)} The candidate sequence is symmetric and summable, so we may form
\[
\sdf(f)=\sum_{\tau\in\Ints}\acvs_\tau e^{-2\pi i\tau f}
=\acvs_0+\acvs_1 e^{-2\pi i f}+\acvs_{-1}e^{2\pi i f}
=1+2a\cos(2\pi f).
\]
A sequence is a valid ACVS of some stationary process iff it is non-negative definite, equivalently (in the summable case) iff its Fourier transform satisfies $\sdf(f)\ge0$ for all $f$. As $f$ ranges over $[-\tfrac12,\tfrac12]$, $\cos(2\pi f)$ takes every value in $[-1,1]$, so
\[
\min_f \sdf(f)=1-2\abs a,\qquad \max_f\sdf(f)=1+2\abs a .
\]
The binding constraint is the minimum: $\sdf(f)\ge0$ for all $f$ $\iff$ $1-2\abs a\ge0$ $\iff$ $\abs a\le\tfrac12$. Therefore
\[
\boxed{\;\{\acvs_\tau\}\ \text{is a valid ACVS}\iff a\in[-\tfrac12,\tfrac12].\;}
\]
(This is exactly the lag-$1$ autocorrelation bound of an $\mathrm{MA}(1)$: $\acf_1=a/\acvs_0=a$ must lie in $[-\tfrac12,\tfrac12]$, attained at the boundary by the unit-coefficient $\mathrm{MA}(1)$ $X_t=\eps_t\pm\eps_{t-1}$, which has $\acf_1=\pm\tfrac12$.)

\end{document}
